\documentclass[12pt,a4paper]{article}

\usepackage[margin=2.5cm]{geometry}
\usepackage[utf8]{inputenc}
\usepackage[brazil]{babel}
\usepackage{lmodern}
\usepackage{fancyhdr}
\usepackage{tocloft}
\usepackage{amsmath}
\usepackage{graphicx}
\usepackage{booktabs}
\usepackage{array}
\usepackage{float}
\usepackage{hyperref}
\usepackage{lastpage}
\usepackage{subcaption}
\usepackage{multirow}

\newcommand{\diameter}{\phi}

\graphicspath{{../../../imagens_gerais/} {images/}} 

\pagestyle{fancy}
\fancyhf{} 
\setlength{\headheight}{52pt}
\lhead{
  \includegraphics[height=1.2cm]{lafe_logo_ingles.jpg} \\
  \small
  Doc ID: PLT0079 \quad
  Doc tipo: Relatório de fabricação \quad
}
\rhead{
  \small \textbf{PLT0079 Hollow-core elíptica (4 internos) - PETG translúcido} \\
  Page \thepage\ of \pageref{LastPage}
}
\lfoot{\small \textcopyright\ 2026 LaFE. All rights reserved.}
\rfoot{\small Última data de revisão: 21/08/2026}

\begin{document}

  Descrição da participação dos envolvidos na fabricação desta fibra:
  \begin{itemize}
    \item Ivan Prearo: Fabricação e puxamento
  \end{itemize}

  \section{Geometria e material}
  \label{sec:geometria}

  Comentários gerais sobre esta etapa:
  \begin{itemize}
    \item Pré-forma:
    \begin{itemize}
      \item Material: PETG esverdeado;
      \item Geometria da secção transversal: \textit{Hollow-core} elíptica;
      \item Comprimento total: \(L_T = 132 mm\);
      \item Comprimento útil: \(L = 100 mm\); 
      \item Diâmetro externo: \(\diameter_{out} = 15 mm\)
      \item Espessura do cladding: \(L_c = 3 mm\);
      \item Diâmetro dos capilares internos (eixo menor): \(\diameter_{in} = 3.98 mm\)
      \item Espessura dos capilares internos: \(L_{in} = 0.6 mm\);
      \item Número de capilares internos: \(N_{in} = 3\);
      \item Excentricidade dos capilares internos: \(e_{in} = 2\);
      \item Distância entre capilares internos\footnote{Essa distância é calculada como se os capilares internos fossem cilíndricos (\(e_{in} = 1\)), no entanto a distância real é, em geral, maior do que a escrita aqui. Estout trabalhando para obter cálculos mais precisos.}: \(D_{in} = 0.2 mm\)
      \item Raio do eixo maior de \(r_{in}\equiv e_{in}\cdot\phi_{in}/2=3.98 mm\);
      \item Razões entre medidas iniciais na Tabela \ref{tab:puxamento1_razoes_das_dimensoes}
    \end{itemize}
    \item Objetivo: Fabricação de fibra hollow-core com capilares elípticos e pressurização.
    \item O esquema da geometria da secção transversal e a foto da pré-forma são mostrados em Fig.~\ref{fig:detalhes_secao_transversal}.
  \end{itemize}

  \begin{table}[H]
  \centering
  \begin{tabular}{c|c|c}
    Dimensão & Medida [mm] & Razão com $\diameter_e$ \\
    \hline
    $\diameter_{out}$ & 70 & 1 \\
    $L_c$ & 12 & 0.171 \\
    $\diameter_{in}$ & 9.6 & 0.137 \\
    $L_{in}$ & 0.6 & 0.0086
  \end{tabular}
  \caption{Relações entre medidas relevantes da preforma}
  \label{tab:puxamento1_razoes_das_dimensoes}
  \end{table}

  \begin{figure}[H]
    \centering
    \begin{subfigure}[t]{0.45\linewidth}
      \centering
      
\includegraphics[height=0.9\linewidth]{images/Secao_transversal.png}
      \caption{Seção transversal da fibra.}
      \label{fig:secao_transversal}
    \end{subfigure}
    \begin{subfigure}[t]{0.45\linewidth}
      \centering
      \includegraphics[height=0.9\linewidth]{images/Medidas_anotadas.png}
      \caption{Medidas relevantes da estrutura.}
      \label{fig:medidas_anotadas}
    \end{subfigure}
    \caption{Detalhamento da seção transversal da preforma.}
    \label{fig:detalhes_secao_transversal}
  \end{figure}

  Como um dos objetivos principais desse puxamento é a pressurização da amostra de acordo com a Figura \ref{fig:pressoes}.
  Para isso, a parte da preforma voltada à tampa de posicionamento foi fechada de acordo com a Figura \ref{fig:tampa_preforma}, mantendo um volume aberto entre os capilares internos para que haja pressurização igual entre eles (Figs. \ref{fig:vistas} e \ref{fig:vistas_com_anotacoes}).

  \begin{figure}[H]
    \centering
    \begin{subfigure}[t]{0.48\linewidth}
      \centering
      \includegraphics[width=\linewidth]{images/Pressurizacao_anotada.png}
      \caption{Seções de pressurização da fibra.}
      \label{fig:pressoes}
    \end{subfigure}
    \begin{subfigure}[t]{0.48\linewidth}
      \centering
      
\includegraphics[width=\linewidth]{images/Secao_pressurizacao.png}
      \caption{Tampa da preforma.}
      \label{fig:tampa_preforma}
    \end{subfigure}
    
    \begin{subfigure}[t]{0.49\linewidth}
      \centering
      \includegraphics[height=0.9\linewidth]{images/Vistas.png}
      \caption{Vistas relevantes da preforma.}
      \label{fig:vistas}
    \end{subfigure}
    \begin{subfigure}[t]{0.49\linewidth}
      \centering
      \includegraphics[height=0.9\linewidth]{images/Vistas_anotadas.png}
      \caption{Fig. \ref{fig:vistas} com anotações.}
      \label{fig:vistas_com_anotacoes}
    \end{subfigure}
    \caption{Detalhamento da seção transversal da preforma.}
    \label{fig:detalhamento_pressurizacao}
  \end{figure}


  \section{Puxamento}
  \label{sec:puxamento}

  Comentários gerais sobre esse puxamento:
  \begin{itemize}
    \item Data de puxamento: 30/07/2026.
    \item A rampa de puxamento é mostrada na Tabela \ref{tab:puxamento_rampa}.
    \item O puxamento resultou em sucesso, produzindo dezenas de metros de fibra a diferentes condições de pressurização antes do rompimento.
  \end{itemize}

  O puxamento ocorreu normalmente até a estabilização razoável do diâmetro da fibra enquanto puxada no \textit{capstan}, que inclui a banda lixo e de transição (Fig. \ref{fig:plot_diametro}).
  Após a estabilidade no diâmetro, foram fabricadas três bandas com pressões incrementais, sendo 2, 3 e 4 mBar para as bandas A, B e C, respectivamente (Fig. \ref{fig:plot_pressao}).
  Durante a fabricação da banda C, próximo ao alvo de pressurização, houve o rompimento da fibra.

  \begin{table}[H]
    \centering
    \begin{tabular}{ccc}   
      \toprule 
      \textbf{Fase} & \textbf{[°C]} & \textbf{Tempo [min]} \\ 
      \midrule
      \multirow{2}{*}{Secagem} & 60 & 60 \\
                                & 80 & 20\\
      \midrule
      \multirow{2}{*}{Escoamento} & 100 & 5 \\
                                & 110 & 2 \\ 
      \midrule
      \multirow{1}{*}{Puxamento} & 120 & --- \\
      \bottomrule
    \end{tabular}
    \caption{Rampa de temperatura no forno da torre antes e durante o puxamento.}
    \label{tab:puxamento_rampa}
  \end{table}

  \begin{figure}[H]
    \centering
    \includegraphics[width=\linewidth]{Graficos/Diametro.png}
    \caption{Diâmetro da fibra ao longo do puxamento.}
    \label{fig:plot_diametro}
  \end{figure}

  \begin{figure}[H]
    \centering
    \includegraphics[width=\linewidth]{Graficos/Pressao.png}
    \caption{Pressão na preforma ao longo do puxamento.}
    \label{fig:plot_pressao}
  \end{figure}

  \section{Resultados}
  \label{sec:puxamento_resultados}

  Comentários gerais dos resultados:
  \begin{itemize}
    \item As medidas absolutas e relativas relevantes estão na Tabela \ref{tab:medidas_abs_rel}.
    \item As micrografias das fibras produzidas estão na Figura \ref{fig:micrografias}.
  \end{itemize}

  É bastante notável que, de acordo com a Tabela \ref{tab:medidas_abs_rel}, a excentricidade dos capilares internos aumentam consistentemente com a pressão.
  No entanto, é necessário tomar tal informação com cautela, uma vez que a excentricidade dos capilares de uma mesma micrografia variam consideravelmente. No caso da banda C (Fig. \ref{fig:micrografia_C}), a menor excentricidade encontrada foi 1.34, enquanto a maior foi 1.68.
  Mesmo com a excentricidade variando para valores inferiores aos sem aplicação de pressão, a média das excentricidades aumentar é significativo.

  Também é importante manter em mente que os valores de excentricidade são calculados considerando uma elípse (como na Fig. \ref{fig:medidas_anotadas}), mas algumas micrografias mostram os capilares internos formando formatos mais circulares (Fig. \ref{fig:micrografia_B}), devido a tensão superficial do PETG.

  \begin{table}[H]
    \centering
    \begin{tabular}{c|c|cc|cc|cc|cc}
      \multirow{2}{*}{Medida} & Inicial & \multicolumn{2}{c|}{Transição} & \multicolumn{2}{c|}{Banda A} & \multicolumn{2}{c|}{Banda B} & \multicolumn{2}{c}{Banda C} \\
      \cline{2-10}
      & Rel. & Abs. & Rel. & Abs. & Rel. & Abs. & Rel. & Abs. & Rel. \\
      \hline
      $\phi_{out}$       & 1.000
                   & 596 & 1.000
                   & 364 & 1.000
                   & 369 & 1.000
                   & 331 & 1.000 \\
      $r_{in}$        & 0.265
               & 122 & 0.205
               & 78.8 & 0.216
               & 76.6 & 0.208
               & 77.4 & 0.234 \\
      $\phi_{in}$       & 0.265
                  & 176 & 0.296
                  & 109 & 0.298
                  & 104 & 0.282
                  & 101 & 0.306 \\
      $L_{in}$        & 0.040
               & 54.8 & 0.090
               & 30.0 & 0.083
               & 34.9 & 0.095
               & 26.3 & 0.079 \\
      $L_{c}$        & 0.200
              & 125 & 0.210
              & 72.2 & 0.198
              & 71.7 & 0.194
              & 60.4 & 0.182 \\
      $e_{in}$       & 2.00
              & 1.38 & ---
              & 1.45 & ---
              & 1.49 & ---
              & 1.52 & --- \\
    \end{tabular}
    \caption{Medidas absolutas e relativas de cada banda para todas as grandezas relevantes.
    A medida de excentricidade não contém unidades, enquanto as outras medidas absolutas estão em $\mu$m.}
    \label{tab:medidas_abs_rel}
  \end{table}

  \begin{figure}[H]
    \centering
    \begin{subfigure}[t]{0.48\linewidth}
      \centering
      \includegraphics[width=\linewidth]{Fotos/Banda_transicao.png}
      \caption{Banda de transição.}
      \label{fig:micrografia_transicao}
    \end{subfigure}
    \begin{subfigure}[t]{0.48\linewidth}
      \centering
      \includegraphics[width=\linewidth]{Fotos/Banda_A.png}
      \caption{Banda A.}
      \label{fig:micrografia_A}
    \end{subfigure}
    \begin{subfigure}[t]{0.48\linewidth}
      \centering
      \includegraphics[width=\linewidth]{Fotos/Banda_B.png}
      \caption{Banda B.}
      \label{fig:micrografia_B}
    \end{subfigure}
    \begin{subfigure}[t]{0.48\linewidth}
      \centering
      \includegraphics[width=\linewidth]{Fotos/Banda_C.png}
      \caption{Banda C.}
      \label{fig:micrografia_C}
    \end{subfigure}
    \caption{Micrografias do final de cada banda.}
    \label{fig:micrografias}
  \end{figure}

\end{document}